\section{Conclusiones}
El modelo identificado a partir de las mediciones entre J5 y J7 proporcionó la
base común para calcular el controlador, los observadores y las simulaciones.
Los mismos coeficientes se programaron en el PSoC para las cuatro
configuraciones. Los registros de 256 muestras mantuvieron un período de
\SI{2}{\milli\second}, no presentaron banderas y contaron con disparo físico y
medición simultánea de J5 y J7.

La comparación directa utilizó una referencia ideal con estados verdaderos y
mostró la dinámica de $A-BK$ sin intervención de un observador. Una vez que el
error de estimación se hace despreciable, las respuestas comparten dicha
dinámica; las diferencias sostenidas entre simulación y banco se explican por
la cuantización, el ruido y la calibración. Los valores finales experimentales
fueron 46,12; 49,62; 46,24 y \SI{50.30}{\milli\volt}, con picos entre 60 y
65 mV.

La evolución de los estados estimados mostró que $\hat x_2$ acompaña la salida
medida, mientras que $\hat x_1$ representa el estado interno no medido. El error
observable RMS descendió de 2,91 a \SI{2.29}{\milli\volt} entre Predictor y
Actual E1, y de 2,91 a \SI{1.77}{\milli\volt} en E2. Ante la misma perturbación
inicial, el uso de la muestra presente redujo inmediatamente la norma del error:
11,69 frente a 22,36 mV en E1 y 8,37 frente a 22,36 mV en E2. Después de la
convergencia no corresponde esperar dinámicas distintas, pues domina $A-BK$.

La estimación modal al 2\,\% produjo tres muestras, o 6 ms, para E1 y E2. Con
el criterio persistente $\|e\|_2<10^{-4}$ V, la perturbación considerada
convergió en 6 ms para E1 y 8 ms para E2. Esta diferencia muestra que el cruce
de un umbral finito también depende de la dirección del error inicial y de la
combinación modal.

Los mandos aplicados tuvieron valores RMS entre 49,78 y 53,28 mV. En las
últimas 50 muestras, la desviación estándar del error observable fue 3,38;
2,49; 3,49 y 1,72 mV para Predictor E1, Actual E1, Predictor E2 y Actual E2.
Estos valores describen las condiciones de los ensayos realizados y muestran
una menor dispersión para Actual E2. El peor tiempo de CPU fue 112 $\mu$s,
equivalente al 5,60\,\% de $T_s$, con un margen temporal mínimo de 94,40\,\%.

La resolución del actuador explica parte de la semejanza entre las respuestas
experimentales. El nivel aplicado coincidió entre predictor y actual en
46,48\,\% de las muestras E1 y 61,33\,\% de las E2; entre solicitudes distintas,
la compresión al mismo código alcanzó 37,73\,\% y 54,17\,\%. Los pasos de
10,56--15,37 mV son grandes frente a la referencia de
\SI{50}{\milli\volt}, de modo que diferencias de mando pueden resultar
invisibles en la salida física.
